\section{Introduction}
\label{sec:introduction}

Financial markets are often studied as complex systems in which heterogeneous agents, institutional constraints, liquidity cycles, information flows and feedback mechanisms interact across multiple time scales \cite{mantegna_stanley_1999_book,cont_2001_empirical}. Even when the analysis is restricted to daily asset returns, market variation also has a relational component: how assets move together across the cross-section, summarized by a correlation matrix. Empirical correlations can contain broad co-movement, localized organization and sampling variation; Random Matrix Theory (RMT) provides a benchmark for separating components compatible with random estimation effects from components that merit structural interpretation \cite{laloux_1999_noise,plerou_2002_random,bun_2017_cleaning}.

The distinction matters because financial risk is inherently multivariate. The volatility of an individual asset matters, but portfolio risk, contagion, diversification and systemic vulnerability depend on how assets co-move \cite{billio_2012_econometric}. During calm periods, correlations may appear moderate and dispersed. During stress periods, correlations often increase, reducing diversification benefits and lowering the effective dimensionality of the market \cite{junior_2012_correlation,zheng_2012_changes}.

Brazil is a useful setting because B3 (Brasil, Bolsa, Balc\~ao), the country's exchange and financial-market infrastructure provider\footnote{\url{https://www.b3.com.br/}}, contains firms from commodity-linked, financial, utility, industrial, consumer, real-estate and telecommunications activities. This diversity makes it possible to ask whether a broad market component coexists with localized dependency patterns. Brazilian studies have examined interbank networks \cite{tabak_2010_topology}, multiscale correlations \cite{pereira_2019_multiscale}, agrarian cross-correlations \cite{stosic_2015_correlations} and equity-market topology \cite{leonel_2022_brazilian}. The specific gap addressed here is narrower: to our knowledge, few B3 equity studies jointly combine RMT filtering, planar financial networks and a clearly qualified set of retrospective model assessments.

Correlation matrices are noisy, especially when the number of assets is not negligible relative to the number of observations. Random Matrix Theory (RMT) offers a benchmark for this problem by comparing the eigenspectrum with the spectrum expected from random correlation matrices \cite{marcenko_1967_distribution,laloux_1999_noise,plerou_2002_random}. Eigenvalues lying inside the Mar\v{c}enko--Pastur band are compatible with random sampling noise under the benchmark model, whereas eigenvalues outside the band indicate candidate collective modes that deserve further interpretation. Subsequent work has extended RMT-based methods for cleaning large covariance matrices \cite{bun_2017_cleaning}. We use RMT as an intermediate layer for reconstructing broad market, group-level and noise-compatible components before graph construction, not as a mechanical classifier of economic sectors.

Network filtering gives an additional representation of the same dependency system. By transforming correlations into distances, financial networks retain selected dependency links among assets and make it possible to inspect market backbones, local clusters and central nodes. The Minimum Spanning Tree (MST) emphasizes the sparsest connected structure \cite{mantegna_1999_hierarchical}, while the Planar Maximally Filtered Graph (PMFG) retains a denser planar representation with richer local connectivity \cite{tumminello_2005_tool,tumminello_2010_correlation}. We use these constructions to ask whether strong dependencies are organized by sectors, whether central nodes change after market-mode removal and whether the intermediate-scale organization of the market contains information beyond average correlation. Filtered networks are also relevant for risk analysis, where peripheral assets can offer better diversification properties than central hubs \cite{pozzi_2013_spread,aste_2010_correlation}.

We develop an integrated econophysics workflow for Brazilian equities traded on B3. It documents stylized facts and correlation structure; applies RMT, clustering and filtered networks to interpret the dependency matrix; and then conducts chronological retrospective model assessments. The structural descriptors in those experiments are estimated from the full synchronized sample. Consequently, the experiments ask whether those descriptors are associated with model performance, not whether they deliver information available at a historical forecast origin.


The paper makes four contributions:
\begin{enumerate}
  \item We build a reproducible econophysics workflow for B3 equities using adjusted daily prices with raw historical coverage from 1998 to 2025, including return construction, stylized-fact checks, correlation analysis, sectoral comparison and rolling synchronization.
  \item We apply Random Matrix Theory to the complete 58-asset synchronized panel (23 August 2006--15 August 2024) and reconstruct market, group/sector and noise-compatible components of the empirical correlation matrix, yielding a spectral interpretation of Brazilian equity dependencies.
  \item We compare MST and PMFG networks built from original and RMT-filtered matrices, with emphasis on hub reorganization, clique structure, clustering, sectoral retention and subsector-level aggregation.
  \item We conduct a chronological, nested-feature-set retrospective assessment of whether full-sample RMT and network descriptors are associated with realized-volatility model performance. This motivates, but does not replace, a point-in-time forecasting design with recursively estimated features.
\end{enumerate}

The article is organized as follows. Section~\ref{sec:theoretical_background} establishes the theoretical background and conceptual foundations. Section~\ref{sec:related_work} reviews the relevant literature. Section~\ref{sec:methodology} presents the proposed reproducible integrated framework, including the data and asset universe, spectral decomposition, filtered networks and retrospective model assessments. The subsequent sections report the empirical results for stylized facts, RMT, clustering, network topology, subsector aggregation and volatility models, followed by discussion, limitations and conclusion.
